% unit U7
\section{Front Ends}\label{sec:frontends}
\mathplain

Report~1 ended with the remark that a compiler could emit evidence for the subsumption steps it takes and have it checked. The front ends are such a compiler, for a model of the Scala typer. Each tree has one, a Lean library that turns a program in this report's notation into a derivation of its calculus and runs it on the machine of \DOTMNF. Its theorems apply theorems of the calculus to the derivation it returns or speak of its own functions, and none adds a fact about the calculus.

A program is a Lean term inside a quotation, \lean{cap\%}, \lean{cc\%} or \lean{clsProg\%}. The front end of \treeCC reads the escape of Section~\ref{sec:cc-later} as \lean{EscSrc}, where \lean{File} abbreviates $\mu(s.\,\{\mathit{read} : \capt{(\forall(u : \top)\,\top)}{\cset{s}}\})$:
\begin{lstlisting}[mathescape=true]
cc% λ(g : ⊤).
  let cb : ($\forall$(f : File ^ {any}) ($\forall$(u : ⊤) File ^ {any}) ^ {any}) ^ {}
         = λ(f : File ^ {any}). λ(u : ⊤). f
  in cb
\end{lstlisting}
In a quotation \lstinline|T ^ {C}| is $T^{C}$. \lean{cb} is the callback \lstinline{f => () => f}, and $\lambda(g : \top)$ stands for the scope that calls \lstinline{withFile}. The $\cany$ of the parameter of \lean{cb} reads as the callback's arrow binder $\kappa$. The $\cany$s of its result read as $\varrho_g$, the scope binder of the body of $\lambda(g : \top)$ (Section~\ref{sec:levels-rule}). The callback has the type $\forall(\kappa, f : \capt{\mathit{File}}{\cset{\kappa}})\,\capt{(\forall(\kappa', u : \top)\,\capt{\mathit{File}}{\cset{f}})}{\cset{f}}$, and the front end binds \lean{cb} at it. Widening \lean{cb} to the type the $\kw{let}$ writes needs $\cset{f} <:_{c} \cset{\varrho_g}$, which the front end rejects with a certificate (Section~\ref{sec:fe-reject}).

\subsection{The pipeline}\label{sec:fe-pipeline}

A program of \treeCL starts with a header that declares the classifiers, the platform, and optionally a use set and a kind. Name resolution inserts a $\kw{let}$ for every operand that is not a variable. It leaves a partial term $p$ whose empty slots are the annotations the program did not write: lambda domains, self shapes and field types. The typer types a term whose slots are all written. The elaborator fills the empty slots from expected types, types each candidate with the typer, and returns the first. That is an elaborated term $\mathit{tm}$ and a derivation of \DOTMNF about $t_c$, the term $\mathit{tm}$ without its $\kw{let}$ types, ascriptions and self shapes. The fill $\bar{p}$ is $p$ with its empty slots read off $\mathit{tm}$.

The typer follows the subtype checker, the level test (in \treeCC and \treeCL), avoidance and box inference (Section~\ref{sec:fe-input}) of the Scala 3 compiler at commit \lean{4dae25087d} \cite{scala3cc}. It tries alternatives in the compiler's order, and no theorem relates the two.

\subsection{Soundness is the result type}\label{sec:fe-sound}

The function \lean{compile} returns a derivation of the calculus, not a yes or no. So whatever \lean{compile} returns is typed by construction. A bug in the typer can make it reject a typable program or run out of fuel, never return an untyped one. The price is that every guarantee is about the term $t_c$ of the derivation, which Section~\ref{sec:fe-input} relates to the written program. \lean{compile} returns the fill $\bar{p}$ and a record
\[
\langle \mathit{tm},\ U_c,\ T_c,\ d : \dtyuj{U_c}{\platctx}{t_c}{T_c},\ \ldots \rangle .
\]
Section~\ref{sec:fe-input} gives its last field.

The theorems above the middle rule of Table~\ref{tab:pipeline} take the hypothesis $h$ that \lean{compile} returned the record, and no proof reads $h$. Each is a theorem of the calculus applied to $d$, so it holds for any record. For instance, \lean{compile\_effect\_safety} is the tree's copy of Theorem~\ref{thm:source-effect} at $d$ and at a run from $t_c$. It asks for a platform capability $\kappa$ that is not an element of the list $U_c$. In \treeCL it asks in addition that $U_c$ write no filter, so that $\kappa \notin U_c$ means that $\kappa$ is not a root of $U_c$ (Section~\ref{sec:source-thms}).

\begin{table}[tb]
\footnotesize
\caption{What each pipeline theorem guarantees, and what it does not. The theorems above the middle rule hold for any record \lean{compile} could return. Those below it are statements about \lean{compile}.}
\label{tab:pipeline}
\centering
\setlength{\tabcolsep}{3pt}
\begin{tabular}{@{}p{0.33\textwidth}p{0.42\textwidth}p{0.21\textwidth}@{}}
\toprule
theorem & guarantees & does not guarantee \tabularnewline
\midrule
\raggedright \lean{compile\_checks}, \lean{compile\_erase}, \lean{compile\_uses\_checks} & \raggedright the checker of \FCdot accepts the translation and its use-set evidence, and the translation erases to the record's term & \raggedright that this term is the written program \tabularnewline
\raggedright \lean{compile\_safe}, \lean{compile\_not\_stuck}, \lean{compile\_run\_progress} & \raggedright every reached source state is final or steps (Theorem~\ref{thm:source-safety}) & \raggedright termination, or a run of the translation by \lean{fcStep?} (\S\ref{sec:mech-checks}) \tabularnewline
\raggedright \lean{compile\_capture\_prediction} & \raggedright the matched target use set is covered by the translated $U_c$ (Theorem~\ref{thm:source-pred}) & \raggedright that $U_c$ is least \tabularnewline
\raggedright \lean{compile\_effect\_safety} & \raggedright a platform capability outside $U_c$ is no root of an inspected variable (Theorem~\ref{thm:source-effect}) & \raggedright anything when that variable has no root \tabularnewline
\raggedright \lean{compile\_lvl\_safety} (\treeCC, \treeCL) & \raggedright at each member-free premise $C <:_{c} D$ of $d$, an atom that confines $D$ (\S\ref{sec:fe-reject}) confines $C$ (Theorem~\ref{thm:lvl-source}) & \raggedright anything about premises that use members \tabularnewline
\raggedright \lean{compile\_classified\_effect\_safety} (\treeCL) & \raggedright every root of an inspected variable has a classifier in $\chi$ (Theorem~\ref{thm:cl-source}) & \raggedright that the header's kind was read \tabularnewline
\midrule
\raggedright \lean{compile\_faithful}, \lean{compile\_full}, \lean{compile\_faithful\_written} & \raggedright the fill agrees with the written slots, the record is the elaborator's first candidate, or the typer's synthesis when every slot is written (in \treeCL possibly at the declared use set), the term has the fill's skeleton (\S\ref{sec:fe-input}) & \raggedright that the term means the written program \tabularnewline
\raggedright \lean{elab\_complete\_direct}, \lean{compile\_complete\_direct} & \raggedright a canonical fill (\S\ref{sec:fe-input}) the typer accepts is found from some fuel on & \raggedright completeness for a fill that is not canonical (\lean{Mid}) \tabularnewline
\bottomrule
\end{tabular}
\end{table}

\subsection{What speaks of the input program}\label{sec:fe-input}

Box inference inserts boxes, unboxings and the $\kw{let}$s that bind them, which the program did not write. The skeleton forgets these, along with all annotations, so that the record's term can be compared with the fill. An annotated term is a term of Figure~\ref{fig:source} in which a literal $\nu(x : S.\,d)$ writes its self shape, a $\kw{let}$ may write the type of its variable, and $(t : T)$ ascribes a type. In \treeCC and \treeCL it may also be the unpacking $\klet \langle \kappa, x\rangle = t \kin u$ of Section~\ref{sec:levels-rule}. A superscript $?$ marks a part that may be absent. The skeleton drops every type, shape and set and keeps the binding structure and the labels:
\begingroup
\color{srccol}
\allowdisplaybreaks
\begin{alignat*}{3}
& \mathit{skel}(x) &\quad& = x \\
& \mathit{skel}(\lambda(x : T).\,t) &\quad& = \lambda x.\,\mathit{skel}(t) \\
& \mathit{skel}(\nu(x : \capt{S}{U^{?}}.\,d)) &\quad& = \nu(x.\,\mathit{skel}(d)) \\
& \mathit{skel}(x\,y) &\quad& = x\,y \\
& \mathit{skel}(x.a) &\quad& = x.a \\
& \mathit{skel}(\klet x\,(: T)^{?} = t \kin u) &\quad& = \mathit{skel}(u)[x \mapsto y] &\qquad& \text{if } \mathit{skel}(t) = y \\
& \mathit{skel}(\klet x\,(: T)^{?} = t \kin u) &\quad& = \klet x = \mathit{skel}(t) \kin \mathit{skel}(u) &\qquad& \text{otherwise} \\
& \mathit{skel}(\klet \langle \kappa, x\rangle = t \kin u) &\quad& = \mathit{skel}(\klet x = t \kin u) &\qquad& \text{in \treeCC and \treeCL} \\
& \mathit{skel}(\dbox{x}) &\quad& = x \\
& \mathit{skel}(\dunbox{C^{?}}{x}) &\quad& = x \\
& \mathit{skel}((t : T)) &\quad& = \mathit{skel}(t) \\
& \mathit{skel}(\ddeftyp{A}{S}) &\quad& = \{A\} \\
& \mathit{skel}(\ddeftyp{C}{D}) &\quad& = \{C\} \\
& \mathit{skel}(\ddeftrm{a}{t}) &\quad& = \{a = \mathit{skel}(t)\} \\
& \mathit{skel}(d \wedge d') &\quad& = \mathit{skel}(d) \wedge \mathit{skel}(d')
\end{alignat*}
\endgroup

\lean{compile\_faithful} relates the record to the written program. The program resolves to $p$, and the fill $\bar{p}$ agrees with every slot $p$ writes. The elaborator's first candidate on $p$ has the term, type and use set of the record, except that in \treeCL the record may hold the use set the header declares. The record's last field says that $\mathit{skel}(\mathit{tm}) = \mathit{skel}(\bar{p})$. So its term may differ from the fill only where the skeleton forgets, as where box inference inserts a box. \lean{C7scala\_type} (\treeK, decided) types C7 written with boxed field types but no box term and no unboxing, and the kernel checks that the elaborated term boxes each field and unboxes before the call.

Inference is complete at direct sites, such as call arguments, ascriptions and bodies of $\kw{let}$s. A fill is canonical when each slot the program leaves empty sits at such a site and holds what the elaborator reads there. At a call argument or a field of a literal it must also pass the typer's check, and in \treeK that check must return the fill itself first. Let the program resolve to $p$, and let its fill be canonical, agree with every slot $p$ writes, and be accepted by the typer. In \treeK let the fill also be $p$ with its empty slots read off the typer's term. Then from some fuel on the program compiles as its fill does (\lean{elab\_complete\_direct}, \lean{compile\_complete\_direct}). In \lean{Mid} a lambda is passed to a function whose type is an intersection of two function types. The typer types the written program at the first. With the lambda's domain erased, the elaborator checks the argument at the second and needs a middle type that the program does not write. So the fill of \lean{Mid} is not canonical (\lean{canonAt}, decided), and the erased program is rejected at every budget (\lean{Mid\_rejected}, decided, in all three trees).

\subsection{Rejections and certificates}\label{sec:fe-reject}

A rejection is a statement about an algorithm. The typer runs on a fuel budget that all its goals share, and a run that exhausts it is reported as a recursion limit, not as a rejection. \lean{Alg} is an inductive judgment with one rule for each alternative of the typer's subtype checker. A run that does not exhaust the fuel answers every goal that \lean{Alg} derives (\lean{shape?\_complete} and its kin). So a rejection without exhaustion means that the failing goal has no \lean{Alg} derivation (the \lean{\_reject} theorems). The proviso stays, since an alternative that never ends, tried first, would exhaust every budget. \lean{Alg} is sound for the calculus (\lean{Alg.answer} in \treeK, \lean{Alg.sound} in \treeCC and \treeCL) and not complete for it. A variant of report~1's bad-bounds lambda, $\src{\lambda(x : \dtyp{A}{\top}{\bot}).\ \klet y : \dtyp{B}{\dfld{a}{\top}}{\dfld{a}{\top}} = x \kin y}$, types through the middle type $\src{x.A}$, which it does not write (built by hand). The typer of \treeK rejects it without exhausting the fuel, at a goal with no \lean{Alg} derivation (\lean{E1\_verdict}, \lean{E1\_not\_alg}, decided).

In \treeCC and \treeCL a rejection names a reason, and a level escape comes with a certificate. An atom $r$ confines a source set $D$ in a source context $\Gamma$ when $\lvlle{\trans{\Gamma}}{\capsr{\trans{\Gamma}}{m}{\trans{D}}}{r}$ for every $m$. Let the typer fail at a goal $C <:_{c} D$ in $\Gamma$, where the context is well formed and every atom of $D$ resolves to itself. Then \lean{certify?} looks for an atom $r$ that confines $D$ and for a depth $n$ with $\neg\,\lvlle{\trans{\Gamma}}{\capsr{\trans{\Gamma}}{n}{\trans{C}}}{r}$. If it finds both, the contrapositive of Theorem~\ref{thm:lvl-source} is the certificate: no member-free derivation proves $C <:_{c} D$ in that context (\lean{escape\_rejected\_at} and \lean{certify?\_isSome}, in \treeCC). The certificate speaks only of the goal and context the typer reached.

For the escape, the goal's context is the body of $\lambda(g : \top)$, with scope binder $\varrho_g$, then \lean{cb} and the callback's scope $\varrho_{\mathit{cb}}, \kappa, f : \capt{\mathit{File}}{\cset{\kappa}}$, which is $\Gamma_{\!X4}$ of Example~\ref{ex:lvl-escape} with $\mathit{File}$ for $\top$. The set $\cset{f}$ resolves to $\cset{\kappa}$ at every depth (\lean{Esc\_caps}, built by hand), and $\kappa$ has the level $\varrho_{\mathit{cb}}$, inside $\varrho_g$, while $\cset{\varrho_g}$ resolves to itself. So $r = \varrho_g$ confines $\cset{\varrho_g}$ and not $\cset{f}$, and no member-free derivation proves $\cset{f} <:_{c} \cset{\varrho_g}$ (\lean{Esc\_rejected'}, built by hand from decided facts). \lean{Esc\_compile\_rejected} says, by \lean{rfl}, that \lean{compile} at the default budget returns this certificate, and \lean{Esc\_rejected}, decided, that no budget compiles it.

\subsection{What is not proved}\label{sec:fe-open}

No theorem says that the typer finds the least use set the rules allow, as the READMEs of the \treeK and \treeCL front ends claim. No theorem gives let insertion or box inference a meaning. The classified theorems take the kind as an argument of \lean{compileKinded} or \lean{compileFiltered}, and nothing reads the kind a header of \treeCL declares. Two front-end instances are not known to be informative (Example~\ref{ex:mini-prog2}). ReadK2 of \treeCC binds closures at $\cset{\kappa_1}$ and $\cset{\kappa_2}$ and calls the second, and CE1b of \treeCL is a front-end version of E1r (Example~\ref{ex:cl-try}). Both are proved from decided premises (\lean{ReadK2\_never\_reads\_k1}, \lean{CE1b\_reads\_only\_control}), but no Lean statement gives their inspected variable a root. The safety theorems hold over a platform, not in an open context.
